\section{线程束调度与延迟隐藏}

\subsection{线程束调度器选择可执行的线程束}

流式多处理器同时维护多个驻留线程束。在线程束已有驻留状态之间切换发射机会，可近似视为\term{零开销线程束调度}{zero-overhead warp scheduling}：当前线程束无法继续时，调度器可以迅速选择另一个已经就绪的驻留线程束，无需执行传统 CPU 线程上下文切换所需的完整保存与恢复。

线程束只有在下一条指令具备执行条件时才能成为候选对象，其中一个关键条件是所需操作数已经可用，即操作数已“ready for consumption”。

\begin{mechanismbox}
每个驻留线程已经在流式多处理器的寄存器文件中占有相应状态；线程束还具有程序计数器等控制状态。在线程束之间切换发射机会时，不需要把整个线程上下文先写回内存再恢复。调度器主要是在已有驻留状态之间选择下一批要推进的线程束。
\end{mechanismbox}

CUDA 不向程序员暴露具体的线程束优先级策略。程序可以依赖的是：满足执行条件的驻留线程束会被硬件调度以持续推进，而不能编写依赖某个固定线程束优先顺序的程序。

\subsection{长延迟操作如何造成线程束停顿}

算术指令所需操作数通常要先进入寄存器。若线程束执行高延迟的全局内存加载，而后续指令立即依赖加载结果，那么在数据返回之前，该线程束无法继续执行依赖指令，因而处于等待状态。

类似地，某些需要多个周期的执行资源也可能造成依赖停顿。关键并不在于停顿来源是内存还是计算，而在于：\textbf{当前线程束的下一条指令暂时无法发射。}

\subsection{用其他线程束覆盖等待时间}

\term{延迟隐藏}{latency hiding}的核心是把一个线程束的等待区间与其他线程束的有用执行重叠。设线程束 \(W_0\) 发出高延迟操作后必须等待；只要同一流式多处理器中存在另一个已经就绪的线程束 \(W_1\)，调度器就可以让 \(W_1\) 获得执行机会。若 \(W_1\) 之后也停顿，则继续寻找 \(W_2\)、\(W_3\) 等可执行线程束。

\begin{figure}[H]
  \centering
  \includegraphics[width=0.88\textwidth]{latency_hiding_p25.png}
  \caption{某个线程束遇到高延迟操作后，处理块可以选择其他就绪线程束继续执行，从而把等待时间与有用工作重叠}
\end{figure}

延迟隐藏不缩短一次全局内存访问本身需要的时间；它通过在这段时间内执行其他独立工作来减少执行资源的空闲时间。因而判断它是否有效时，应问“等待期间有没有别的线程束可以推进”，而不是只问“单次操作的延迟是多少”。

\subsection{为什么驻留线程数远大于物理核心数}

一个流式多处理器通常支持的最大驻留线程数明显大于其物理执行核心数，这使硬件能够保持较大的候选工作池。

假设流式多处理器中只有很少线程束。如果这些线程束同时发出长延迟加载，调度器就可能找不到任何就绪线程束，执行资源只能空等。若驻留线程束数量更多，则在部分线程束等待时，其他线程束已经准备好操作数的概率通常更高。

因此可以把关系概括为
\[
\text{更多驻留线程束}
\Longrightarrow
\text{更大的就绪候选集合}
\Longrightarrow
\text{更强的延迟隐藏机会}.
\]

这条关系解释了为什么占用率常常重要：占用率提高意味着流式多处理器上同时驻留的线程束更多，从而更可能在某个线程束停顿时找到替代工作。

\subsection{延迟隐藏失效的典型情形}

延迟隐藏只有在存在其他可执行线程束时才有作用。以下情况会使“增加线程束”无法继续解决问题：

\begin{itemize}
  \item \textbf{内存带宽已经饱和。} 更多线程束只会增加等待请求，不能提高已经受带宽上限约束的吞吐；
  \item \textbf{所有线程束都受到同类未解决依赖阻塞。} 如果没有任何线程束的下一条指令已经就绪，调度器仍然无事可做；
  \item \textbf{缺少足够的指令级并行。} \term{指令级并行}{instruction-level parallelism}不足时，各线程束可能都很快进入依赖停顿，无法形成可交错执行的独立工作。
\end{itemize}

\begin{keybox}
延迟隐藏的必要条件不是“高占用率”这个数字本身，而是驻留工作中持续存在可发射的独立线程束。高占用率通常提高这一机会，但不能越过带宽、依赖或执行资源的硬瓶颈。
\end{keybox}
